\scanpage{50}
\uline{Cor} Tout \deleted{\unsure{conti|cont}} segment saturé contient un onglet -- donc si $\mathcal{R}$ n'est pas la relation grossière il existe au moins deux onglets (car il existe alors au moins deux segments saturés, complémentaires l'un de l'autre)

Soit \inserted{$\unsure{\Sigma_{1}|\Sigma_{0}}$ l'ens des onglets de $R$,} $\deleted{\unsure{S_{0}|S_1}}\unsure{T_{1}|T_{0}}$ la réunion des onglets, $\deleted{\unsure{S_1}}R_1$ son complémentaire. Alors, si $\deleted{S_1}R_1\ne\phi$, $\deleted{S_1}R_1$ a au moins deux segments composants ($\operatorname{card}\unsure{\pi_0|\Pi_0}(R_1)\geq 2$), sinon $R_1$ serait un \deleted{\unsure{s}}segment donc contiendrait un onglet, absurde. Si $\operatorname{card} R_1=2$, alors $R_1\in\mathcal{R}$ \deleted{est un} donc $\mathcal{R}$ consiste en les onglets et en $R_1$, (car \deleted{\unsure{$R_1$|$R$}} les $\unsure{S_i|S'}\subset R_1$ sont de card. $\geq 2$, \deleted{\unsure{puis|puisq}} car ceux de cardinal $1$ seraient des onglets \ldots)

Si $\operatorname{card} R_1\geq 3$, $R_1$ est lui $=$ un graphe polygonal; et $\mathcal{R}_1=\mathcal{R}\vert R_1$ est une relation d'équivalence non panachée. De plus, \deleted{\unsure{al|Au}}Aucun $\unsure{S_i|S_1}\in\mathcal{R}_1$ n'est contenu dans un segment composant de $R_1$ -- sinon il serait contenu dans un segment saturé contenu dans $\deleted{\unsure{S_1}}R_1$ (\unsure{Savoir|savoir} son cocomposant contenu dans ce composant), lequel contiendrait un onglet $\subset R_1$ -- absurde. Si $\mathcal{R}_1$ est la relation grossière, on a fini -- $\mathcal{R}$ est formée des onglets, et de $R_1$.
% margin note written sideways
\margin{\unsure{Construction|Conclusion} \unsure{fausse}}
% archive sheet number: 49

